\question{5.m1}{Voor een willekeurige normale verdeling geldt dat de volgende drie grootheden aan elkaar gelijk moeten zijn:}
\begin{enumerate}[label=(\alph*)]
    \item gemiddelde, standaarddeviatie en mediaan.
    \item modus, de totale kans en de variantie.
    \item mediaan, gemiddelde en modus.
    \item standaarddeviatie, variantie en spreidingsbreedte.
\end{enumerate}
\answer{
    De normale verdeling is een symmetrisch continue kansverdeling. 
    Hierdoor zijn de mediaan, het gemiddelde en de modus allemaal gelijk aan elkaar.       
    Het juiste antwoord is dus (c).
}
